\begin{annexitem}{annex:p3.adm.10}{thm:p03-adm-exterior}{Removing the exterior}
\begin{proof}
For each formula \(\chi\) and its finite free-variable set, a state
specifies which of those variables are assigned inside \(U\), and a
partition of the remaining variables into blocks having equal values
outside \(U\). Different blocks have different values. There are only
finitely many states for a finite free-variable set. We define a formula
\(\chi_t\) in the base signature for every state \(t\). Its free variables
are among the variables marked inside. Its intended truth is that of
\(\chi\) under any assignment realizing the state, with the indicated
inside values in \(M\).

For an atomic formula \(U(x)\), use truth or falsehood according to
whether \(x\) is inside. For \(R(x_1,\ldots,x_n)\), use the same atom
if all variables are inside, and falsehood otherwise. For equality,
retain the equality if both variables are inside, use falsehood if
exactly one is inside, and use truth or falsehood according to whether
the two outside variables belong to the same block. Closed logical
truth and falsehood can be written as \(\forall z\,(z=z)\) and
\(\exists z\,(z\ne z)\). These bound variables do not add a free
variable or a non-logical symbol.

Negation commutes with translation. For a conjunction or disjunction,
restrict \(t\) to each constituent's free-variable set, translate that
constituent in the restricted state, and keep the connective.
For \(\exists x\,\eta\), extend \(t\) in the following finitely many
ways: put \(x\) inside, put it in any one of the already represented
outside blocks, or put it in one new outside block. The translation is
the disjunction of the inside case \(\exists x\,\eta_{t+x\text{ inside}}\)
and the formulas for all the outside cases. In an outside case \(x\)
is no longer a free variable of the translated formula. For
\(\forall x\,\eta\), take instead the conjunction of the inside case
\(\forall x\,\eta_{t+x\text{ inside}}\) and those same outside cases.
If there are no old outside blocks, the fresh outside case is still
present. The newly bound variable is absent from the parent state's
free-variable set, so these extensions have no conflicting assignment.

Induction on \(\chi\) verifies the intended truth equivalence. The
atomic and Boolean cases are exactly the interpretations just specified.
Every possible value of a quantified variable is either inside, equal
to one of the finitely many represented outside values, or different
from all of them. Conversely all these cases can be realized. The last
one can always be realized because the exterior is infinite. By the
induction hypothesis all its realizations have the same truth value
for the matrix, since they have the same state. This proves both the
existential and universal steps, independently of the cardinality of
the infinite exterior.

The construction stays in \(A\). The finite set of states of a finite
variable set is obtained by finite combinatorics; blocks can be coded
by finite sequences of natural numbers and canonically renumbered when
a state is restricted. At each subformula node, compute its whole finite
family of translated formulas from the already computed child families.
The formation operations and state restrictions have \(\Sigma\) graphs.
Set recursion and syntactic closure from
\cref{thm:p03-adm-recursion,thm:p03-adm-syntax} therefore put these
families in \(A\), including at an infinitary connective whose
constituent set belongs to \(A\). If the inside universe is empty, its existential cases are false
and its universal cases true, exactly as required; the exterior remains
infinite. For a sentence the free-variable set is empty and has one state. Its translation is the required
\(\theta^\circ\). No exterior predicate or new relation symbol remains.
\end{proof}
\end{annexitem}

\begin{annexitem}{annex:p3.adm.11}{thm:p03-adm-relativized-separation}{Relativized separation}
\begin{proof}
Let the projective presentations be \(\phi_1,\phi_2\), with respective
expanded signatures and designated domains. Rename auxiliary symbols
so the two expansions are independent. Work in one new ambient universe,
with a shared predicate \(U\) for the common reduct and separate private
predicates \(D_1,D_2\) for the two expansion universes.

For each side also make private copies of all its relation symbols,
including every base relation \(R\in\sigma\). This detail matters:
the original sentence \(\phi_i\) can constrain a base symbol on tuples
outside its designated reduct. Those values must not be forced to be
false or identified with the other expansion's values. Keep the original
symbols \(R\in\sigma\) as separate shared predicates for the reduct.
The designated-domain predicate of each \(\phi_i\) is read as the
shared \(U\).

Let \(\operatorname{Good}_i\) be the following sentence in this ambient
language. The domain \(U\) is contained in \(D_i\); either may be empty.
The privately renamed \(\phi_i\), with all its quantifiers relativized
to \(D_i\), holds. Finally, for each \(R\in\sigma\), the private copy
\(R_i\) agrees with shared \(R\) on tuples entirely in \(U\):
\[
 \forall\bar x\,
 \left(\bigwedge_j U(x_j)\ \longrightarrow\
                   (R_i(\bar x)\leftrightarrow R(\bar x))\right).
\]
Let \(\operatorname{Pure}\) say that each shared \(R\) is false on
tuples not entirely in \(U\). These are sentences in \(L_A\): renaming
and relativization are structural recursions on formula codes, and
replacement collects the displayed clauses for the signature set
\(\sigma\in A\). Their conjunction has no free variables.

The implication
\[
 \operatorname{Good}_1\land\operatorname{Pure}
 \quad\longrightarrow\quad
 \neg(\operatorname{Good}_2\land\operatorname{Pure})
\]
is valid. Otherwise the two induced structures on \(D_1,D_2\), using
their private predicates, would witness membership of the same induced
structure on \(U\) in both disjoint classes. Relativization preserves
truth by induction on the source formulas, including their infinitary
connectives. Empty induced domains have the source's empty-structure
semantics; they have not been discarded.

Apply \cref{thm:p03-adm-interpolation} to the ordinary nonempty
ambient universes. The predicates for induced domains need not be
nonempty. All private symbols on the two
sides are disjoint, so the interpolant \(\theta\in A\) uses only
\(U\) and shared symbols from \(\sigma\). At this point it can still
depend on the ambient universe outside \(U\); ordinary interpolation
alone does not yet finish the proof.

Apply \cref{thm:p03-adm-exterior} to obtain \(\theta^\circ\in L_A(\sigma)\).
If \(M\) belongs to the first class, choose its witnessing expansion,
identify the designated reduct with \(M\), and add a disjoint infinite
set of extra objects to the ambient universe. Interpret \(D_1\) as
the original expansion universe, its private relations as in that
expansion, and shared relations as those of \(M\), false outside
\(U=M\). Interpret unused symbols arbitrarily. This structure satisfies
\(\operatorname{Good}_1\land\operatorname{Pure}\), hence \(\theta\).
Its reduct to the interpolant's language has a pure infinite exterior,
so \(M\models\theta^\circ\). For a member of the second class, perform
the same construction with its witnessing expansion and \(D_2\).
It satisfies \(\operatorname{Good}_2\land\operatorname{Pure}\), so the
second interpolating implication forces \(\neg\theta\), and therefore
\(M\not\models\theta^\circ\). This proves separation. If the two classes
are complementary, the separator defines the first one, which is the
Souslin--Kleene assertion.

For function and constant signatures first use the relational graph
translation of \cref{def:p02-w02c-p02-003-relational-image} on each
expansion and on the base signature. Require the usual function-graph
and singleton-constant axioms on each relevant domain. The two resulting
projective classes of relational structures remain disjoint and are
separated by the result just proved. Inverse graph translation of the
separator defines the desired class of original structures, by
\cref{thm:p02-w02c-p02-003-translations}; both translations preserve
membership in \(A\). Finally a sentence using a subset of a finite
base signature has finitely many non-logical symbols, so the finite-symbol
restriction needed in Barwise's application is retained.
\end{proof}
\end{annexitem}
